%% AAMAS 2027 main-track draft, story-led version.
%% Layout uses acmart sigconf (the class the official AAMAS template is built on), so page
%% counts match. To submit: move the body into the official aamas_2027 template
%% (\documentclass[sigconf,anonymous]{aamas}).
%% Figures in figures/ are exported from the interactive explainer (light theme, vector PDF).
\documentclass[sigconf,anonymous,acmthm=false]{acmart}

\settopmatter{printacmref=false}
\setcopyright{none}
\renewcommand\footnotetextcopyrightpermission[1]{}
\acmConference[AAMAS '27]{Proc.\ of the 26th International Conference on Autonomous Agents and Multiagent Systems (AAMAS 2027)}{May 3--7, 2027}{Hanoi, Vietnam}
\acmYear{2027}

\usepackage{booktabs}
\usepackage{graphicx}
\usepackage{subcaption}
\usepackage[most]{tcolorbox}
\usepackage[capitalise,nameinlink]{cleveref}

%% ---------- Theorem environments ----------
\theoremstyle{acmplain}
\newtheorem{theorem}{Theorem}
\newtheorem{lemma}{Lemma}
\newtheorem{proposition}{Proposition}
\newtheorem{corollary}{Corollary}
\theoremstyle{acmdefinition}
\newtheorem{definition}{Definition}
\emergencystretch=1.5em
\raggedbottom

%% ---------- Story beats ----------
\definecolor{storyblue}{HTML}{2448C7}
\definecolor{storybg}{HTML}{EEF2FC}
\newtcolorbox{story}[1][]{enhanced, breakable, colback=storybg, colframe=storyblue,
  boxrule=0pt, leftrule=1.6pt, sharp corners, left=5pt, right=5pt, top=3pt, bottom=3pt,
  before skip=6pt, after skip=6pt, fontupper=\small, #1}
\newcommand{\plain}[1]{\par\smallskip\noindent\textit{In plain words:} #1\par\smallskip}

%% ---------- Notation ----------
\newcommand{\E}{\mathbb{E}}
\newcommand{\Pb}{\mathbb{P}}
\newcommand{\TV}{\mathrm{TV}}
\newcommand{\dH}{d_{\mathrm{H}}}
\newcommand{\Kflat}{K_{\mathrm{flat}}}
\newcommand{\dmin}{d_{\min}}
\newcommand{\shat}{\hat{s}}
\newcommand{\cube}{\{0,1\}}
\newcommand{\A}{\textsf{Alpha}}
\newcommand{\B}{\textsf{Beta}}
\newcommand{\G}{\textsf{Gamma}}

\begin{document}

\title{How Much Memory Does Accountability Cost?}
\subtitle{Three Reviewers, One Sticky Note, and the Price of Compressed Multi-Agent Records}

\begin{abstract}
When a team of AI agents compacts its execution history to fit a context window, the
compacted record becomes the only evidence left for deciding, after the fact, who did
the work. Contract theory designs payments for a given signal, and trace compression
optimises task performance with no incentive term, so no existing account says what a
record must keep for payments read from it to reward effort. Here, we follow three
reviewer agents and a coordinator whose memory holds one of $K$ labels, and we model the
record as a stochastic encoder of the team's effort profile scored on two jobs at once:
reconstruction error and the \emph{margin}, the extra pay an agent is guaranteed for
working, whatever its teammates do. By eliminating the payment rule, we show that the
margin equals the least total-variation shift that an agent's lone deviation causes in the
record, which turns every design question into an optimisation over the record alone.
With noisy outcomes, the squared margins of all agents sum to at most a constant times
the record's entropy, so accountability costs a fixed number of bits per agent; every
record that is symmetric in the agents supports no budget-balanced payment at any size;
and with exact outcomes, accountability becomes free at exactly the 2-domination number
of the hypercube (4, 10 and 32 labels for 3, 5 and 7 agents). On reviews written by a
language model, deleting 30 bytes of authorship labels raises attribution error by
0.943 bits (95\% CI $[0.802, 1.089]$). Together, the results show that accountability is
priced in agent identity, not in content volume.
\end{abstract}

\keywords{Moral hazard in teams; contract design; budget balance; information structures; accountability; compressed records; LLM multi-agent systems; domination number}

\maketitle

%% =====================================================================
\section{Introduction}\label{sec:intro}
%% =====================================================================

\begin{story}
\textbf{Monday, 09:02.} A coordinator agent receives a pull request and splits the review
three ways. \A{} checks the API changes, \B{} the tests, \G{} the security-sensitive
diff. Each sub-agent works through tens of thousands of tokens and returns a verdict. By
09:40 the coordinator's context window is nearly full, so it does what every long-running
agent does: it replaces the three traces with a short summary of themselves, a sticky
note. A week later the change causes an incident. Someone asks which reviewer actually
read the diff, whose verdicts should be trusted next time, and whom the team's
reputation system should reward. The sticky note is all that is left.
\end{story}

Multi-agent systems built from language models routinely discard their own execution
histories: sub-agents return summaries rather than traces, and coordinators compact older
exchanges when their context nears its limit. Unlike a storage bill, a context window
cannot be extended, so the budget on what survives is hard. The surviving record then has
two jobs (\cref{fig:pipeline}). The first is \emph{reconstruction}: recover what the team
did accurately enough to act on. The second is \emph{accountability}: keep enough about
each agent that a later, bounded payment, read from the record alone, still rewards that
agent's effort whatever its teammates did. The two jobs pull in opposite directions. A
reconstruction-optimal record spends its few labels well by merging profiles that are
cheap to confuse; accountability needs each agent, acting alone, to visibly move the
record.

\begin{figure*}[t]
\centering
\includegraphics[width=0.86\textwidth]{figures/pipeline.pdf}
\caption{The story's object of study. The coordinator dispatches three reviewers, then
compacts their traces into a record $Y$ holding one of $K$ labels ($\log_2 K$ bits).
Everything to the left of the dashed line is gone after compaction. Two readers later use
the same label: one reconstructs what the team did, the other pays each agent.}
\label{fig:pipeline}
\end{figure*}

Contract theory has long studied how to reward unobserved effort. Moral hazard in
teams~\cite{holmstrom1982teams} shows that a joint output hides individual effort, and
algorithmic contract
theory~\cite{babaioff2006combinatorial,castiglioni2023multiagent,dutting2025multiagent,dutting2025combinatorial,goel2024weighted,dutting2024survey}
computes optimal payments for a given outcome technology. The economics of information
ranks signals by their value~\cite{holmstrom1979moral,blackwell1953comparisons}. In each
case the monitoring signal is taken as given, so none of these lines asks which signal a
system can afford to keep under a hard size budget.

Work on agent systems approaches the record from the other side. Failure attribution in
LLM multi-agent systems~\cite{zhang2025which,zhang2025agentracer,cemri2025fail} recovers
who caused a failure from a trace, and full traces improve attribution by up to 76.5\%
over partial ones~\cite{chen2026elephant}. Trace compression~\cite{hong2026trace,lin2026tracer}
shortens traces to preserve task performance. Both assume that whatever the record keeps
is chosen for its usefulness to the task; neither includes an incentive term. The gap
already causes failures in deployed systems: in Bittensor, validators who copied a
merged consensus were reported to out-earn honest ones, and the fix that was deployed
changed what the record kept rather than how closely validators were
monitored~\cite{opentensor2024weightcopy,liu2024bittensor,opentensor2025commit}.

The question this leaves open is the one the coordinator faces: \emph{given a budget on
the record, which records can be kept, and which of them still support payment?}
Answering it is harder than either line of work alone, because the record and the
payment rule are both unknown and interact multiplicatively; in our experiments a direct
search over both did not terminate at seven agents.

In this work, we model the record as a stochastic encoder that maps the team's effort
profile to one of $K$ labels, and we score it on reconstruction error and on the
\emph{accountability margin}: the extra pay an agent is guaranteed for working rather than
shirking, whatever its teammates do, under some bounded payment rule that reads only the
record. Following the three reviewers through six episodes, we establish:
\begin{enumerate}
\item \textbf{You are paid only for what the record sees} (\cref{thm:elim}). The payment
rule can be eliminated: an agent's margin equals the least total-variation shift its lone
deviation causes in the record, against the most confusing mixture of teammate behaviour.
\item \textbf{One bit costs accuracy at a constant rate} (\cref{thm:onebit}). With one
retained bit the reconstruction-accountability frontier is linear and in closed form, and
no margin above $1/n$ is possible.
\item \textbf{When accountability is free} (\cref{thm:flat,thm:prior}). With exact outcomes
the least record size that is reconstruction-optimal and still pays is the 2-domination
number $\gamma_2(Q_n)$ of the hypercube, with certificates at $n=3,5,7$; we characterise
the threshold for every full-support prior.
\item \textbf{Luck makes memory linear} (\cref{thm:noisy}). With noisy outcomes the squared
margins of all agents sum to at most a constant times the record's entropy, so each agent
held to account costs a constant number of bits and every sublinear summary eventually
pays for nothing.
\item \textbf{A fixed pot needs names} (\cref{thm:bb}). Under budget balance the margin is
bounded by how differently the record reacts to deviations by different agents, so every
record symmetric in the agents fails at any size.
\item \textbf{Measured records.} Common compaction mechanisms are easy to make
unaccountable by accident, and on reviews written by a language model, deleting 30 bytes
of authorship labels raises attribution error by 0.943 bits.
\end{enumerate}
Every number is recomputed by an independent implementation (135 checks), and 22
predictions were registered and hashed before each harness was written. The rest of this
paper is organised as follows. \Cref{sec:model} formalises the story and defines the
margin. \Cref{sec:elim,sec:onebit,sec:flat,sec:noisy,sec:bb} present the five theoretical
results, \cref{sec:experiments} measures deployed-style records, and
\cref{sec:related,sec:design} discuss related work and design rules.

%% =====================================================================
\section{The Story, Formalised}\label{sec:model}
%% =====================================================================

\begin{story}
The coordinator's note can show one of a handful of labels. Later, a reputation module
reads the same note and pays each reviewer between 0 and 1 credit. Each reviewer knows
this in advance. Reading the diff properly costs a reviewer real compute; returning a
plausible verdict without reading it costs almost nothing. A reviewer reads the diff only
if doing so is guaranteed to earn at least that much extra credit.
\end{story}

\paragraph{Effort and records.} A team of $n$ agents produces an \emph{effort profile}
$S\in\cube^n$, where $S_i=1$ means agent $i$ did its work; we write profiles with \A, \B,
\G{} as agents $1,2,3$ from left to right, so $011$ means \A{} shirked and the others
worked. $S$ is uniform on $\cube^n$ unless stated (\cref{sec:prior} treats other priors).
A \emph{record} of size $K$ is a stochastic encoder $q:\cube^n\to\Delta([K])$ with
$q[s,y]=\Pb(Y=y\mid S=s)$. The label $Y$ is everything that persists: the encoder's coins
are not retained, and any seed or identifier kept outside the record counts toward
$K$.\footnote{Otherwise one private bit per agent stored beside a single shared label would
satisfy every condition below at $K=1$.} A fixed-length encoding costs $\log_2 K$ bits.

\paragraph{Reconstruction.} A rule $\shat:[K]\to\cube^n$ has expected Hamming error
$D(q,\shat)=\E[\dH(S,\shat(Y))]$, and $D_0(K)=\min_{q,\shat}D(q,\shat)$ is the classical
$K$-median value of the cube.

\paragraph{Payment and the margin.} Agent $i$ is paid by a bounded decoder
$p_i\in[0,1]^K$ that also sees only $Y$. The \emph{deviation vector} of agent $i$ at
teammate profile $s_{-i}$,
\[ d_i(s_{-i})[y]=q[(1,s_{-i}),y]-q[(0,s_{-i}),y], \]
is the change in the label distribution caused by $i$ alone. The record \emph{implements
margin} $\mu$ if bounded decoders exist with
\begin{equation}\label{eq:M}\tag{M}
\langle p_i, d_i(s_{-i})\rangle \ge \mu \quad\text{for every } i \text{ and every } s_{-i}.
\end{equation}
The left side is the expected pay of $i$ when working minus when shirking, with
teammates fixed at $s_{-i}$. \plain{the margin is the extra pay an agent is
\emph{guaranteed} for working instead of shirking, whatever the others do. It is built in
three steps: compare the two worlds for each teammate situation, keep the smallest gap,
and choose the payment rule that makes that smallest gap largest. Because pay is capped
at 1, the margin lies in $[0,1]$.}

\Cref{fig:alpha}(c) and \cref{tab:alpha} work this out for \A{} under a record that keeps
only the number of reviewers who worked. The best rule pays $1/3$ credit per worker, and
every one of \A's four teammate situations gives a gap of exactly $1/3$, so \A's margin
is $1/3$.

\begin{table}[t]
\caption{\A's four worlds under the success-count record ($n=3$), with the best bounded
payment (1/3 credit per recorded worker). The margin is the smallest gap.}
\label{tab:alpha}
\footnotesize
\begin{tabular}{@{}lccc@{}}
\toprule
teammates (\B, \G) & pay if \A{} works & pay if \A{} shirks & gap\\
\midrule
both shirk & 0.33 & 0.00 & 0.33\\
only \B{} works & 0.67 & 0.33 & 0.33\\
only \G{} works & 0.67 & 0.33 & 0.33\\
both work & 1.00 & 0.67 & 0.33\\
\bottomrule
\end{tabular}
\end{table}

\begin{definition}[Settlement game]\label{def:game}
Fix a record $q$. Each agent $i$ chooses effort $a_i\in\cube$ at cost $c_i\ge0$; the
profile draws $Y\sim q[a,\cdot]$; agent $i$ is paid $p_i(Y)\in[0,1]$ and gets
$\E[p_i(Y)]-c_ia_i$. A payment rule \emph{ex-post implements} effort by $i$ if working is
a best response whatever the teammates do: $\langle p_i,d_i(s_{-i})\rangle\ge c_i$ for
every $s_{-i}$, which is \eqref{eq:M} at $\mu=c_i$.
\end{definition}

\begin{proposition}[The margin is a price]\label{prop:threshold}
Write $\mu_i(q)$ for the largest margin $q$ implements for agent $i$. Effort costing $c_i$
is ex-post implementable from $q$ by bounded payments if and only if $\mu_i(q)\ge c_i$.
When this holds for every agent, all-work is an ex-post equilibrium in which participation
is automatic, since $\E[p_i\mid a_i=1]\ge c_i+\E[p_i\mid a_i=0]\ge c_i$. No mechanism whose
transfers are functions of $Y$ does better, because every such transfer is a bounded
decoder after rescaling.
\end{proposition}

\plain{$\mu_i(q)$ is a price tag: the most costly piece of work the record can ever pay
for. The cost $c$ is how much more doing the work properly costs the agent than cutting
corners, measured in the same credits as the pay.}

The object of study is the least error at a required margin,
$D^*(K,\mu)=\min\{D(q,\shat):(q,\shat)\text{ implements margin }\mu\}$. A record is
\emph{flat} at $(n,K)$ if it attains $D_0(K)$ while implementing some $\mu>0$, and
$\Kflat(n)$ is the least $K$ at which a flat record exists: the size at which
accountability is free. \Cref{sec:onebit,sec:flat} use the ex-post notion because with
exact outcomes the two-label record ``everyone worked or not'' already makes all-work a
Nash equilibrium with margin 1; \cref{sec:noisy} shows that with noisy outcomes even Nash
accountability needs memory linear in $n$.

%% =====================================================================
\section{You Are Paid Only for What the Record Sees}\label{sec:elim}
%% =====================================================================

\begin{story}
The coordinator could keep a full list of who worked (8 labels) or something smaller.
\Cref{fig:alpha}(a,b) shows two 4-label notes with exactly the same reconstruction error.
Under one, every reviewer can be rewarded; under the other, \A{} cannot be paid for reading
the diff at any price. Nothing about their error reveals which is which.
\end{story}

\begin{figure*}[t]
\centering
\begin{subfigure}[b]{0.31\textwidth}\centering
\includegraphics[width=\linewidth,page=1]{figures/hero_flat.pdf}
\caption{Accountable record: $D=0.5$, $\mu=2/3$.}\end{subfigure}\hfill
\begin{subfigure}[b]{0.31\textwidth}\centering
\includegraphics[width=\linewidth,page=1]{figures/hero_near.pdf}
\caption{Reconstruction-only record: $D=0.5$, $\mu=0$.}\end{subfigure}\hfill
\begin{subfigure}[b]{0.33\textwidth}\centering
\includegraphics[width=\linewidth,page=1]{figures/margin_gaps.pdf}
\caption{\A's gaps under the count record.}\end{subfigure}
\caption{Same size, same error, opposite accountability. Corners are effort profiles
(\A, \B, \G); filled corners are the four labels' reconstruction targets. Edge thickness is
the total variation of the record across that edge; red dashed edges are \emph{dead}: one
agent's change of mind leaves the record's output unchanged. (a) spreads each blank corner
uniformly over its three filled neighbours; (b) sends it to one nearest filled corner.
(c) The gap in each of \A's four teammate situations under the count record
(\cref{tab:alpha}); the smallest gap is the margin, and the dashed line is an effort cost
of 0.25.}
\label{fig:alpha}
\end{figure*}

The design problem couples the record $q$ and the decoders $p_i$, whose product makes a
direct search bilinear. The first result removes $p_i$.

\begin{theorem}[Payment-decoder elimination]\label{thm:elim}
Fix a record $q$ and an agent $i$. A bounded decoder $p_i\in[0,1]^K$ satisfying
\eqref{eq:M} exists if and only if
\[ \min_{\lambda\in\Delta(\cube^{n-1})}\tfrac12\big\|\E_\lambda[d_i]\big\|_1\ge\mu,
\qquad \E_\lambda[d_i]=\textstyle\sum_{s_{-i}}\lambda(s_{-i})\,d_i(s_{-i}). \]
\end{theorem}

\begin{proof}
The inner minimum in $\max_p\min_{s_{-i}}\langle p,d_i(s_{-i})\rangle$ is over a finite
set, so it equals the minimum over its convex hull, i.e.\ over $\lambda$. The map
$(p,\lambda)\mapsto\langle p,\E_\lambda[d_i]\rangle$ is bilinear on the compact convex sets
$[0,1]^K\times\Delta$, so by von Neumann's minimax theorem
$\max_p\min_\lambda\langle p,\E_\lambda[d_i]\rangle=\min_\lambda\sum_y(\E_\lambda[d_i])_y^+$,
since a linear form on the unit cube is maximised by $p_y=\mathbf1[(\E_\lambda d_i)_y>0]$.
Each $d_i(s_{-i})$ is a difference of probability vectors, so its coordinates sum to zero,
as do those of any mixture; for a zero-sum vector $v$, $\sum_yv_y^+=\tfrac12\|v\|_1$.
\end{proof}

\plain{picture, for each teammate situation, the bag of labels the record draws from when
\A{} works and the bag when \A{} shirks. The share of the bag one would have to relabel to
turn one into the other is the total variation. Paying 1 credit exactly on the labels that
are more common when \A{} works converts that share into an equal gap in pay, and the cap
of 1 credit means no rule does better. The adversary then picks the most confusing mixture
of what \B{} and \G{} do. So a record supports accountability exactly to the extent that
each agent, acting alone, visibly moves it.}

\begin{corollary}[Edge condition]\label{cor:edge}
Taking $\lambda$ a point mass, it is necessary that
$\TV\big(q[(1,s_{-i}),\cdot],q[(0,s_{-i}),\cdot]\big)\ge\mu$ for every $i$ and $s_{-i}$. In
particular, a record with a single dead edge in direction $i$ gives agent $i$ margin $0$
under every bounded payment.
\end{corollary}

\Cref{fig:alpha}(b) has three dead edges, so its margin is zero although its error is
optimal. The theorem also explains the count's weakness. One agent always moves a count by
one, which suggests full margin, but the minimising mixture makes the teammate count
uniform on $\{0,\dots,n-1\}$; shifted by one, that distribution overlaps itself except at
the two ends, so the count's margin is exactly $1/n$ (\cref{tab:alpha}). Computationally,
the margin of any record is one small linear program per agent, the instrument used
throughout \cref{sec:experiments}; the reduced design question is a feasibility problem
over $q$ alone and runs in milliseconds.

%% =====================================================================
\section{One Bit: Accuracy for Fairness at a Constant Rate}\label{sec:onebit}
%% =====================================================================

\begin{story}
Suppose the note is a single tick. Majority vote, ``tick if most reviewers worked'', is the
best one-bit summary of the team, but if \B{} and \G{} both worked the tick is there
whatever \A{} does. A random spot-check, ``pick one reviewer and tick if that reviewer
worked'', is fair to everyone and a poor summary of the team.
\end{story}

\begin{theorem}[One-bit frontier]\label{thm:onebit}
Let $n$ be odd, $S$ uniform on $\cube^n$ and $K=2$, with arbitrary bounded payments and
arbitrary reconstruction. Put $M_n=\E[\min(|S|,n-|S|)]$. Then for $0\le\mu\le1/n$,
\[ D^*(2,\mu)=M_n+n\mu\Big(\tfrac{n-1}{2}-M_n\Big), \]
and no record implements margin $\mu>1/n$, whatever the payment rule.
\end{theorem}

\paragraph{Why it is linear.} With $K=2$ the record is one function $q:\cube^n\to[0,1]$
and the margin condition becomes one orientation per agent. Aligning orientations by input
flips turns the feasible set into a shift, by $\mu|s|$, of the cone of monotone functions;
for a fixed reconstruction the error is linear in $q$, so the optimum moves along a line.
The optimal record mixes majority vote and the random-bit sampler in proportion $1-n\mu$
to $n\mu$:
$q^*(s)=\mu|s|+(1-n\mu)\mathbf1\{|s|>n/2\}$ (\cref{fig:onebit}). The full proof fixes the
reconstruction pair before optimising the record and closes with a Rademacher coupling
(supplement, Section A). An exhaustive LP over all reconstruction orbits matches the
closed form to within $3.1\times10^{-15}$ at $n=3,5,7,9$.

\plain{every extra unit of fairness costs the same amount of accuracy, and at the ceiling
$\mu=1/n$ the record beats blind guessing by exactly half a bit \emph{in total}.}

\begin{figure*}[t]
\centering
\begin{subfigure}[b]{0.45\textwidth}\centering
\includegraphics[width=\linewidth,page=1]{figures/onebit_frontier.pdf}
\caption{Exact frontier $D^*(2,\mu)$ for $n=3,5,7,9$.}\end{subfigure}\hfill
\begin{subfigure}[b]{0.45\textwidth}\centering
\includegraphics[width=\linewidth,page=1]{figures/onebit_record.pdf}
\caption{The optimal record at $n=5$, $\mu=0.08$.}\end{subfigure}
\caption{The one-bit frontier (\cref{thm:onebit}). (a) Each line ends at $(1/n,(n-1)/2)$;
no record implements a larger margin. The open circle marks $\mu=0.08$ at $n=5$, where the
least error rises from $M_5=1.5625$ to $1.7375$. (b) Probability that the bit is 1, by the
number of agents who worked: the random-bit sampler (blue, weight $n\mu$) adds exactly $\mu$
per worker, and majority vote (green, weight $1-n\mu$) supplies the rest.}
\label{fig:onebit}
\end{figure*}

%% =====================================================================
\section{When Is Accountability Free?}\label{sec:flat}
%% =====================================================================

\begin{story}
Give the coordinator more labels. Each label now stands for one best-guess profile, a
\emph{centre}, and every profile is sent to a nearest centre. Can the note be as accurate as
any $K$-label note and still pay every reviewer?
\end{story}

\begin{lemma}[Necessity]\label{lem:nec}
Let $D_0$ be strictly decreasing at $K$, so the $K$ centres are distinct. If some
non-centre $s$ has a unique nearest centre, the record is not flat, for any bounded payment.
\end{lemma}
\begin{proof}
Step from $s$ one bit towards its unique nearest centre $c$ to get $s'$; $c$ is still the
unique nearest centre of $s'$. Attaining $D_0(K)$ forces $q[s,\cdot]$ and $q[s',\cdot]$ to
be the point mass on $c$'s label, so that edge has $d_i=0$ and \cref{cor:edge} gives margin 0.
\end{proof}

When the centres $C$ dominate the cube, \cref{lem:nec} requires every vertex outside $C$
to have at least \emph{two} neighbours in $C$: $C$ is a \emph{2-dominating set} of $Q_n$.
The converse holds constructively.

\begin{lemma}[Sufficiency]\label{lem:suff}
Let $C$ be 2-dominating with $|C|=K$. The record that keeps each centre's own label and
spreads every other profile uniformly over the labels of its neighbouring centres, decoded
to the centre and paying agent $i$ the $i$-th bit of the centre, has $D=D_0(K)=(2^n-K)/2^n$
and margin $1-1/\dmin\ge1/2$, where $\dmin=\min_{t\notin C}|N(t)\cap C|$.
\end{lemma}

\begin{proof}[Proof idea]
Centres are reconstructed exactly and every other profile costs one bit, which meets the
counting bound $D_0(K)\ge(2^n-K)/2^n$. Along an edge in direction $i$ between two
non-centres, every centre-neighbour of the working end has bit $i$ equal to 1 and every one
of the shirking end has it equal to 0, so the payment jumps by 1; when one end is a centre
the other puts weight $1/d$ on it and the jump is $1-1/d$ (supplement, Section A).
\end{proof}

\plain{treat profiles as houses, edges as roads and centres as fire stations. A house with a
single nearest station is a trap: stepping towards the station leaves the label unchanged,
so the agent who stepped is invisible. The fix is a backup: every house without a station
must be next door to at least two stations, with a coin flip between them.}

\begin{theorem}[Flatness threshold]\label{thm:flat}
Restricted to $K\ge\gamma(Q_n)$, the regime in which some optimal $K$-median code is
dominating, $\Kflat(n)=\gamma_2(Q_n)$, attained with margin $1-1/\dmin\ge1/2$ by the record
of \cref{lem:suff}.
\end{theorem}

Below $\gamma(Q_n)$, small $K$ is excluded by counting: in a flat record every profile at
distance $r$ from $C$ needs two centres at distance $r$, so shell $r$ holds at most
$\lfloor K\binom nr/2\rfloor$ profiles, and filling shells from the inside lower-bounds
$2^nD_0(K)$. This reduces to Fink and Jacobson's bound
$\gamma_k(G)\ge k|V|/(\Delta+k)$~\cite{fink1985ndomination} and is tight at all three
known flat points. With $\gamma(Q_n)$ known for $n\le9$~\cite{ostergard2001length9}, the
theorem holds unconditionally there: $\Kflat(3)=4$, $\Kflat(5)=10$ (with $K=9$ proved
infeasible), $\Kflat(7)=32$ ($K=29,30,31$ proved infeasible), while $\Kflat(9)\in[97,128]$
is open. Flatness saves surprisingly little: at $n=20$ a flat record still keeps about 17 of
20 bits.

\begin{figure*}[t]
\centering
\begin{subfigure}[b]{0.31\textwidth}\centering
\includegraphics[width=\linewidth,page=1]{figures/stations_trap.pdf}
\caption{A trap: optimal, $\mu=0$.}\end{subfigure}\hfill
\begin{subfigure}[b]{0.31\textwidth}\centering
\includegraphics[width=\linewidth,page=1]{figures/stations_cube.pdf}
\caption{2-dominating: flat, $\mu=2/3$.}\end{subfigure}\hfill
\begin{subfigure}[b]{0.31\textwidth}\centering
\includegraphics[width=\linewidth,page=1]{figures/prior_cube.pdf}
\caption{Tilted prior: flat, $\mu=1/3$.}\end{subfigure}
\caption{Fire stations on the 3-cube; filled corners are centres. (a) The optimal code
$\{000,100,010,111\}$ leaves $001$ (red ring) with a unique nearest centre; stepping
towards it gives a dead edge. Numbers count neighbouring centres. (b) The odd-weight class
$\{100,010,001,111\}$: every other corner has three neighbouring centres, so
$\Kflat(3)=\gamma_2(Q_3)=4$. (c) Under i.i.d.\ Bernoulli$(0.7)$ effort the unique optimal
4-code leaves $000$ at distance 2 (label d2$\cdot$3: three nearest centres), yet the record is
flat (\cref{thm:prior}).}
\label{fig:stations}
\end{figure*}

\subsection{Other priors}\label{sec:prior}
Condition \eqref{eq:M} never mentions the prior; only $D$ and $D_0$ do, and
\cref{lem:nec} survives verbatim for any full-support prior $\pi$ at whatever radius the
nearest centre lies. Write $\nu(s)$ for the number of centres nearest to $s\notin C$.

\begin{theorem}[The threshold for every full-support prior]\label{thm:prior}
$\Kflat(\pi)$ is the least $K$ for which some $\pi$-optimal $K$-median code has
$\nu(s)\ge2$ for every $s\notin C$, and spreading each profile uniformly over its nearest
centres attains it.
\end{theorem}

The proof shows that the spread record's margin is $1-\max\nu(s')/\nu(s)$ over edges on
which $s'$ is one step closer to $C$, and that the condition forces $\nu$ to strictly
decrease along such edges (supplement, Lemmas S2 and S3). Tilting can move the threshold
either way (\cref{fig:stations}(c)): at $n=4$, $\Pr(S_i{=}1)=0.7$ lowers it from 7 to 5,
while $0.6$ raises it to 13. We checked both lemmas against brute force on all 46,592
optimal codes at $n=3,4$ under uniform, product and random priors, with no exceptions.

%% =====================================================================
\section{Luck Makes Memory Linear}\label{sec:noisy}
%% =====================================================================

\begin{story}
Real reviews are noisy. A reviewer who reads the diff usually catches the bug ($h=0.9$),
and one who skims sometimes gets lucky ($l=0.1$). The note can record only outcomes. With
100 reviewers and a note that counts successes, one shirker lowers the expected count by
0.8, while luck alone moves it by about $\pm3$: a whisper in a crowded room
(\cref{fig:noise}(a)).
\end{story}

Agent $i$ chooses $a_i$; its recorded outcome $S_i$ is Bernoulli$(h)$ if $a_i=1$ and
Bernoulli$(l)$ if $a_i=0$, independently, with $0\le l<h<1$. Actions reach the label
through $Q[a,\cdot]=\sum_s\Pr(S=s\mid a)\,q[s,\cdot]$, and \cref{thm:elim} applies to $Q$
unchanged. The \emph{Nash margin}
$\mu_i^N=\TV(Q[\mathbf1,\cdot],Q[(0,\mathbf1_{-i}),\cdot])$ is the largest cost at which
working is a best response to a working team; the ex-post margin is never larger.

\begin{theorem}[Accountability needs linear memory]\label{thm:noisy}
Let $Y$ be the label under all-work and $H(Y)$ its entropy in nats. For every record with
$K$ labels,
\[ \sum_{i=1}^n\big(\mu_i^N\big)^2\le\frac{(h-l)^2}{2h(1-h)}H(Y)\le\frac{(h-l)^2}{2h(1-h)}\ln K. \]
\end{theorem}

\begin{proof}
Let $P_i^v$ be the law of $Y$ given $S_i=v$ under all-work. A deviation by $i$ changes only
the law of $S_i$, so $Q[\mathbf1]=hP_i^1+(1-h)P_i^0$, $Q[(0,\mathbf1_{-i})]=lP_i^1+(1-l)P_i^0$
and $\mu_i^N=(h-l)T_i$ with $T_i=\TV(P_i^1,P_i^0)$. With $\bar P=Q[\mathbf1]$, Pinsker's
inequality~\cite{cover2006} gives
$I(Y;S_i)=h\,\mathrm{KL}(P_i^1\Vert\bar P)+(1-h)\,\mathrm{KL}(P_i^0\Vert\bar P)\ge2h(1-h)T_i^2$,
using $\TV(P_i^1,\bar P)=(1-h)T_i$ and $\TV(P_i^0,\bar P)=hT_i$. Under all-work the $S_i$ are
independent, so $I(Y;S_i)\le I(Y;S_i\mid S_{<i})$ and the chain rule gives
$\sum_iI(Y;S_i)\le I(Y;S)\le H(Y)\le\ln K$.
\end{proof}

\plain{to be visible, an agent must leave a trace in the record; independent agents' traces
cannot share room; and the record holds only so much. So each agent held to account uses up
a fixed slice of the record.}

\begin{corollary}\label{cor:noisy}
(i) All-work is a Nash equilibrium at cost $c$ under some bounded payment reading a
$K$-label record only if $\log_2K\ge n\cdot\frac{2h(1-h)}{\ln2}\big(\frac c{h-l}\big)^2$.
(ii) If $\log K=o(n)$, as for a success count ($\log_2(n+1)$ bits) or any fixed-size
summary, the largest cost the record can sustain for every agent tends to 0 as $n$ grows.
\end{corollary}

This is Holmstr\"om's team problem with the hiding done by compression rather than by the
production technology, and it is quantitative. Block counts attain the order.

\begin{proposition}[Block counts]\label{prop:blocks}
Recording the number of successes in each block of $m$ agents uses
$\log_2K=(n/m)\log_2(m+1)$ bits and gives every agent Nash margin
$(h-l)\max_k\Pr[\mathrm{Bin}(m-1,h)=k]$.
\end{proposition}

As shown in \cref{fig:noise}(b), at $n=100$ one bit per agent keeps the largest possible
margin $h-l=0.8$, blocks of 5 keep $0.5249$ with 51.7 bits, the count keeps $0.1055$ with
6.7 bits, and ``all succeeded'' keeps less than $10^{-4}$. At $n=1000$ the blocks of 5 still
give $0.5249$ while the count falls to $0.0336$. Block counts meet \cref{thm:noisy} up to a
factor logarithmic in $1/c$; on 3000 random stochastic records at $n=3$ to 6 the largest
ratio of margin to bound was 0.764.

\begin{figure*}[t]
\centering
\begin{subfigure}[b]{0.47\textwidth}\centering
\includegraphics[width=\linewidth,page=1]{figures/whisper.pdf}
\caption{Count when everyone works vs.\ one shirker, $n=100$.}\end{subfigure}\hfill
\begin{subfigure}[b]{0.5\textwidth}\centering
\includegraphics[width=\linewidth,page=1]{figures/memory.pdf}
\caption{Nash margin against record size, $n=100$.}\end{subfigure}
\caption{What luck does to accountability ($h=0.9$, $l=0.1$). (a) The success count under
all-work (blue) and with one shirker (orange) differ in total variation by only 10.5\%.
(b) Dots are exact Nash margins of block-count records (\cref{prop:blocks}); the grey curve
bounds every record of that size (\cref{thm:noisy}); the shaded region is infeasible for
every record at effort cost $c=0.3$ (amber line).}
\label{fig:noise}
\end{figure*}

\paragraph{Robustness and uniqueness.} The ex-post margin also decides uniqueness: if every
agent's ex-post margin exceeds $c$, some bounded payment makes working strictly dominant,
so all-work is the unique equilibrium, and if some agent's ex-post margin is below $c$, no
bounded rule makes working dominant for it (\cref{thm:elim} applied to $Q$). The flat
record's two margins coincide ($0.5063$ at $n=5$), so whenever it sustains effort it does so
uniquely, whereas the count sustains Nash up to $0.5249$ but dominance only up to $0.1600$.

%% =====================================================================
\section{A Fixed Pot Needs Names}\label{sec:bb}
%% =====================================================================

\begin{story}
The reputation module has a fixed budget per review: 9 credits, always paid out in full.
To reward \A{} for working it must take credit from \B{} or \G. If the note says only ``2 of
3 read the diff'', it looks the same whichever reviewer skimmed, so any credit taken from
``the shirker'' on that note goes to someone who could have been the shirker
(\cref{fig:pot}(a)).
\end{story}

A payment rule is \emph{budget balanced} if $\sum_ip_i(y)=B$ for every label $y$. Write
$Q_1$ for the label law under all-work, $Q_{0,i}$ when $i$ alone shirks, and
$\bar Q_0=\frac1n\sum_jQ_{0,j}$.

\begin{theorem}[Budget balance needs attribution]\label{thm:bb}
If a budget-balanced rule with $p_i\in[0,1]^K$ makes working a best response for every
agent by margin $t$ at all-work, then
$t\le\frac1n\sum_{i=1}^n\TV(Q_{0,i},\bar Q_0)$.
\end{theorem}

\begin{proof}
Let $d_i=Q_1-Q_{0,i}$ and $\bar d=\frac1n\sum_jd_j$. Since $\sum_ip_i=B\mathbf1$ and
$\langle\mathbf1,\bar d\rangle=0$,
$nt\le\sum_i\langle p_i,d_i\rangle=\sum_i\langle p_i,d_i-\bar d\rangle+B\langle\mathbf1,\bar d\rangle=\sum_i\langle p_i,\bar Q_0-Q_{0,i}\rangle$,
and a $[0,1]$-valued $p_i$ against a zero-sum vector gains at most its total variation.
\end{proof}

\plain{detection asks whether the record moves when $i$ shirks; attribution asks whether it
moves \emph{differently} for different agents. With a fixed pot only attribution pays.}

\begin{corollary}\label{cor:bb}
A record symmetric in the agents, meaning its label law is unchanged when agents are
permuted, admits no budget-balanced rule sustaining any positive effort cost, at any size
$K$. This covers success counts, ``all succeeded'', every function of the multiset of
outcomes, and the multiset of artifacts of \cref{sec:experiments}.
\end{corollary}

This is Holmstr\"om's impossibility~\cite{holmstrom1982teams} with the joint output
replaced by any symmetric record, a quantitative counterpart of the attribution
requirement of~\cite{legrosmatthews1993,rahmanobara2010}. As depicted in
\cref{fig:pot}(b), at $n=5$ the count and the all-succeeded record detect deviations as well
as anything (Nash $0.5249$) and attribute none. Among 1200 random records the largest ratio
of budget-balanced margin to the bound was $1.000$, attained by the flat record at $n=3$.

\begin{figure*}[t]
\centering
\begin{subfigure}[b]{0.34\textwidth}\centering
\includegraphics[width=\linewidth,page=1]{figures/pot.pdf}
\caption{A 9-credit pot read from ``how many''.}\end{subfigure}\hfill
\begin{subfigure}[b]{0.6\textwidth}\centering
\includegraphics[width=\linewidth,page=1]{figures/attribution.pdf}
\caption{Detection against attribution at $n=5$.}\end{subfigure}
\caption{A fixed pot needs names. (a) Each agent's gain from working when the note ``2 of 3''
splits the pot as $1/2/6$ credits (\A/\B/\G) and ``3 of 3'' as $3/3/3$: the gains always sum
to zero, so some agent has no reason to work. (b) For each record, Nash margin (blue),
budget-balanced margin (orange) and the attribution bound of \cref{thm:bb} (green), with
$h=0.9$, $l=0.1$.}
\label{fig:pot}
\end{figure*}

\paragraph{The price in payments and in memory.} Say a label \emph{certifies} agent $i$ if
every profile that can produce it has $S_i=1$. For every record, the least expected payment
making $i$ work (teammates working) is at least $ch/(h-l)$, the full-record cost, with
equality if and only if $c\le(h-l)\Pr[Y\text{ certifies }i\mid S_i=1]$ (supplement,
Proposition S2). A count certifies $i$ only when everyone succeeded, so at $c=0.02$ it costs
1.699 times the full record at $n=100$ and 6.989 times at $n=1000$. Combining this with
\cref{thm:noisy}, a principal who gains $v$ per working agent and pays $\kappa$ per stored
bit should hold nobody to account when $v-ch/(h-l)<\kappa\beta$, with
$\beta=\frac{2h(1-h)}{\ln2}(\frac c{h-l})^2$, and everyone with one bit each when
$v-ch/(h-l)>\kappa$: the designer chooses how many agents to keep accountable, and memory
then scales with that number.

%% =====================================================================
\section{What Real Records Do}\label{sec:experiments}
%% =====================================================================

\begin{story}
No coordinator designs its note this way. It keeps a count, a threshold, the first few
reviewers' reports, or a hash into a fixed number of buckets. We measured these.
\end{story}

\paragraph{Compaction mechanisms.} Every profile is enumerated, so all numbers in
\cref{tab:mech} are exact. No mechanism beats the frontier, and none below
$K=\gamma_2(Q_n)$ is flat. The last two rows use the same ten labels and have the same
error $D_0(10)=0.6875$; one admits margin $1/2$ and the other admits none, for any bounded
decoder. Keeping the first $m$ agents' bits gives margin exactly 0 for every $m<n$, so
prefix-4 is reconstruction-optimal at its budget and worthless for accountability.

\begin{table}[t]
\caption{Compaction mechanisms at $n=5$, uniform prior, no noise. ``Flat'': $D=D_0(K)$
and $\mu>0$.}
\label{tab:mech}
\footnotesize
\begin{tabular}{@{}lrrrrl@{}}
\toprule
mechanism & $K$ & $D$ & $D_0(K)$ & $\mu$ & flat\\
\midrule
full profile & 32 & 0.0000 & 0.0000 & 1.0000 & yes\\
count & 6 & 1.5625 & 0.8750 & 0.2000 & no\\
threshold & 2 & 1.5625 & 1.5625 & 0.0000 & no\\
prefix-4 & 16 & 0.5000 & 0.5000 & 0.0000 & no\\
top-2 & 16 & 0.7188 & 0.5000 & 0.0000 & no\\
hash-10 & 10 & 1.5000 & 0.6875 & 0.0000 & no\\
\midrule
flat (\cref{lem:suff}) & 10 & 0.6875 & 0.6875 & 0.5000 & \textbf{yes}\\
$K$-median & 10 & 0.6875 & 0.6875 & 0.0000 & no\\
\bottomrule
\end{tabular}
\end{table}

\paragraph{Episodes.} To test whether one bit of effort per agent is too coarse, a
coordinator fans out $n$ subtasks, each agent returns an artifact whose quality depends on
its effort, and a record compresses the whole trace. At $n=5$ the full trace and one bit
per agent both have margin $0.5000$, while the multiset of artifacts, which keeps every
artifact's content and loses only authorship, has $0.1425$, the count $0.1190$ and the
final output alone $0.0243$ (\cref{fig:measured}(b)). Across matched-storage comparisons
at $n=3$ to 7 the relative margin lost by removing identity exceeds the relative storage
saved (ratio at least 1.07, median 1.4 to 2.3).

\begin{figure*}[t]
\centering
\begin{subfigure}[b]{0.33\textwidth}\centering
\includegraphics[width=\linewidth,page=1]{figures/mech_scatter.pdf}
\caption{Mechanisms at $n=5$ (\cref{tab:mech}).}\end{subfigure}\hfill
\begin{subfigure}[b]{0.62\textwidth}\centering
\includegraphics[width=\linewidth,page=1]{figures/episode.pdf}
\caption{Episode records at $n=5$: margin and storage.}\end{subfigure}
\caption{Identity, not volume. (a) Reconstruction error against margin; the flat and
$K$-median records share an error of $0.6875$ and differ by the whole margin. (b) The
multiset of artifacts keeps all content and loses only authorship, and keeps less than a
third of the margin of one bit per agent.}
\label{fig:measured}
\end{figure*}

\paragraph{Reviews a model actually wrote.} We closed the loop on the opening story. In a
three-agent code review written by codestral-2508 (192 episodes), the same model at
temperature 0 judged which reviewers read the diff. The full trace and the artifact
multiset hold the same text and differ by about 30 bytes of ``Reviewer $k$:'' labels;
attribution error rises from 0.391 to 1.333 bits, a gap of 0.943 bits (95\% bootstrap CI
$[0.802,1.089]$). The multiset recovers \emph{how many} reviewers worked about as well as
the full trace (0.438 against 0.391 bits); the gap is \emph{which}. It is symmetric in the
agents, so by \cref{cor:bb} no budget-balanced rule can settle from it at any size. A trained
stylometric classifier might attribute better than a general-purpose
judge~\cite{dietrich2026peers}. Of 14 registered conjectures, 11 were falsified, including
our own prediction that the count has full margin; the supplement lists all of them.

%% =====================================================================
\section{Related Work}\label{sec:related}
%% =====================================================================

Holmstr\"om~\cite{holmstrom1982teams} showed that when joint output hides individual
effort, no budget-balanced sharing rule makes effort a best response; here per-agent
information existed and was destroyed by a budget. Georgiadis and
Szentes~\cite{georgiadisszentes2020} let a principal buy signals about one agent at
constant marginal cost; Kadan and Swinkels~\cite{kadanswinkels2008} characterise optimal
contracts under our payment bounds for a given information structure; Legros and
Matthews~\cite{legrosmatthews1993} and Rahman and Obara~\cite{rahmanobara2010} show that
budget-balanced partnerships need attributable deviations, which \cref{thm:bb} prices for
compressed records. The informativeness principle~\cite{holmstrom1979moral} and pairwise
full-rank conditions~\cite{fudenberg1994folk} require that deviations move the signal;
\cref{thm:elim} is a quantitative one-shot version robust to teammate mixtures. Blumrosen,
Nisan and Segal~\cite{blumrosen2007auctions} bound auctions with few-bit messages and Fadel
and Segal~\cite{fadelsegal2009} price the communication that incentives cost for hidden
types; we price the memory that accountability costs for hidden action. Differential
privacy~\cite{dwork2006calibrating} bounds each individual's influence on a record from
above; accountability bounds it from below. Cooperative multi-agent reinforcement learning
decomposes a shared reward~\cite{foerster2018coma,sunehag2018vdn,rashid2018qmix,li2025nucleolus};
there the per-agent signal was never produced. For 2-domination we use known
bounds~\cite{fink1985ndomination,chellali2012survey,bujtas2016bounds} and domination numbers
of small cubes~\cite{ostergard2001length9}, and claim no novelty for the invariant; the
contribution is the reduction to it.

%% =====================================================================
\section{Design Rules and Epilogue}\label{sec:design}
%% =====================================================================

\begin{story}
Back on Monday, the coordinator could have kept three ticks, one per reviewer, at a cost of
three bits. It kept a summary that listed what was found and dropped who found it.
\end{story}

The results give five rules for anyone compacting agent records. \textbf{R1, preserve
identity before content:} dropping authorship alone loses most of the margin and all
budget-balanced settlement (\cref{cor:bb}). \textbf{R2, keep one witness per agent:} one bit
per agent is flat with margin 1, while bits for a fixed subset give margin exactly 0.
\textbf{R3, do not settle from a global count:} its margin is $1/n$ without noise and
vanishes with noise. \textbf{R4, budget memory per accountable agent:} each needs at least
$\beta$ bits (\cref{thm:noisy}). \textbf{R5, measure accountability separately:} equal
reconstruction error can hide margins of $1/2$ and $0$ (\cref{tab:mech}); the linear
program of \cref{thm:elim} takes milliseconds.

\paragraph{Limitations.} Effort is binary and loss is Hamming, although
\cref{thm:elim,thm:bb} hold for any finite outcome space and \cref{thm:noisy} uses binary
outcomes only in its Pinsker step; \cref{thm:onebit} needs the uniform prior. Records are
one-shot and self-contained: repetition, audits or side channels each restore incentives
and change the question, and a strategic compactor is outside the model. The experiments
show such records are easy to build by accident, not that a deployed system has built one.

\paragraph{Conclusion.} A size budget can erase each agent's footprint while leaving
reconstruction optimal. Accountability costs bits per agent, and budget balance costs
identity. Open problems are $\Kflat(9)\in[97,128]$, the one-bit frontier under non-uniform
priors, and the logarithmic gap in \cref{thm:noisy}.

\label{endbody}\bibliographystyle{ACM-Reference-Format}
\bibliography{refs}

\end{document}
